% !TeX root = ../papers/nc-p3.tex
\section{例外交換表示の非正則な八直線}
\label{sec:exceptional-exchange-sheets}
\begin{lemma}\label{lem:bisymmetric-reducible-quick-test}
隣接階数6の両対称反転窓で、いずれかの反対核が二つの一次式の積を含むとする。
標準Koszul左右AS正則となる必要十分条件は、関係空間に純テンソルがなく、
両側の全ての非零評価が階数3以上であることである。
\end{lemma}
\begin{proof}
必要性は反転窓の二つの成長障害による。平方因子の場合、軸と相手の評価階数は
3であり、定理\ref{thm:double-plane-ore}の三変数完全交叉の条件に反すれば
極形式系に階数1が現れ、全窓に純テンソルを与える。
相異なる二因子なら補題\ref{lem:reciprocal-apolar-factors}で四軸の階数は3となる。
二つの直交三平面の交わりが二次元以上なら、二行二列関係の二平面に
行列式零の非零元があり、これも純テンソルとなる。
横断する場合の$\kappa_U=0$、および根基直線の場合の$\det t=0$も、
それぞれ純テンソルを与える。従って定理\ref{thm:transverse-factor-amalgam},
\ref{thm:radical-factor-amalgam}の十分条件が成立する。
\end{proof}

\begin{theorem}\label{thm:exceptional-exchange-eight-lines}
$\lambda=1,A_1=A_0$の全交換表示で、$zD_0(z^2-1)\ne0$とする。
非標準ASとなる$[a:b:c:d]$は下の八直線の和集合との交わりに等しい。
\[
\begin{array}{c|c|c|c}
 &\text{非零座標の位置}&\text{直線の式}& r^2\\\hline
\text{左}&02&a=b,\ c=b&b\\
 &03&a=zb,\ c=z^2b&b\\
 &12&c=b,\ d=b&b\\
 &13&d=zb,\ c=z^2b&b\\\hline
\text{右}&02&a=b,\ d=zb&b\\
 &03&a=zb,\ d=b&b\\
 &12&c=za,\ d=za&za\\
 &13&a=zd,\ c=zd&zd
\end{array}
\]
四つの例外交換表示には同じ判定が前節の変換で適用される。各表示で八直線は
全て重根二次曲面上にあり、特にその表示の単根窓は全て正則である。
\end{theorem}
\begin{proof}
純テンソルがないことは前節の角ブロックの証明で係数差の非零を必要としない。
反対切断の二つの長方形ブロックの各行（転置した側では各列）は零でない。
実際、前節の四行二列表示の第一側は各行に$1$または$z$があり、第二側でも同様である。
いずれかで二行が比例すれば、その二文字の平方の非自明な関係式を生じ、
二つの異なる一次因子を持つ。補題\ref{lem:bisymmetric-reducible-quick-test}が使える。
全ての行対が独立なら両対角鉛筆は滑らかで定理\ref{thm:smooth-apolar-complete}による。
従ってあとは評価核が二次元以上になる条件を決定すればよい。

左評価を前節の$u,v,s,t,L,K$で書く。一方の半分$u,v$が零の場合の排除、
および両半分が非零の場合の核から$s,t$への射影の単射性は二行列の可逆性だけで成立する。
$u_0u_1v_0v_1\ne0$とする。$c_0=u_1/u_0$, $v'=v_1/v_0$と置くと
係数比較から
\[
 c_0 L_{00}=K_{11}v',\quad c_0L_{01}v'=K_{10},\quad
 c_0K_{01}v'=L_{10},\quad c_0K_{00}=L_{11}v' .
\]
これらを掛け合わせて整理すれば、行列の成分が零でも前節と同じ行列式の矛盾になる。
$v_0v_1=0$で$u_0u_1\ne0$なら、例えば$v_1=0$で第一と第三の積式の
$t_0$係数、第ニと第四の積式の$t_0$係数を比較すると
$L_{00}=L_{10}=0$になる。$v_0=0$も同様である。
逆に$u$の一座標が零なら厚い式の対応する二行目または一行目により、
$v$の一座標が零である。右評価は$K$の行と列を同時に交換して同じ結論を与える。
従って評価核の次元超過は、低い一座標を1、高い一座標を$r\ne0$とする
表の八場合に限られる。

この座標形では薄い二式を使い相手の低い二座標を高い二座標から一意に解ける。
残る厚い四式が恒等的に消える条件を取る。例えば左02では
$r^2(a-b)=r^2(a-c)=0$と
$D_0=r^2(d-c)=r^2(d-b)$を得る。$D_0\ne0$なので表の通りとなる。
\path{computations/verify_exceptional_exchange_sheets.py}は八場合全てで、
$r^2$だけを因子とする係数から独立な二つの直線条件を取り出し、
残りが$D_0\ne0$の下で表の$r^2$を必要十分とすることを検算する。
右辺は該当直線の可逆部分で非零であり、幾何点では平方根が二つある。
各直線が前節の八差の平面のうち異なる二枚に含まれ、$G$がその上で消えることも
表に代入すれば分かる。$z\ne0,\pm1$ではこれらの直線は相異なる。
\end{proof}
